\section{Métricas SaaS}

Al constituirse como una plataforma \textit{Software as a Service} (SaaS) con un modelo de negocio de empresa a empresa (B2B), la salud financiera y la viabilidad a largo plazo de SocioUnido no se evalúan bajo las métricas del comercio tradicional, sino a través de los \textit{Unit Economics} propios de los ingresos recurrentes.

En el mercado de SaaS, la retención es tan crítica como la adquisición. Si una plataforma pierde clientes (clubes) más rápido de lo que los adquiere, o si el costo de convencer a una comisión directiva supera el valor que ese club aportará a lo largo del tiempo, el modelo colapsa. Para evitar esto, proyectamos y mediremos nuestro desempeño utilizando los siguientes indicadores clave de rendimiento (KPIs).

\subsection{Ingresos recurrentes y crecimiento (MRR y New Net MRR)}

La métrica principal de estabilidad financiera es el \textbf{\textit{Monthly Recurring Revenue} (MRR)}, que representa la suma de los cargos fijos mensuales cobrados a todas las instituciones:
\begin{equation}
    \text{MRR} = \sum_{i=1}^{n} \text{Cuota Mensual Recurrente}_i
\end{equation}
donde $n$ representa la cantidad total de clubes suscriptos y activos en el mes corriente. El MRR garantiza que las operaciones diarias se ejecuten con previsibilidad financiera sin depender de ingresos transaccionales extraordinarios.

Dado el modelo híbrido de SocioUnido, el crecimiento financiero neto, denominado \textbf{\textit{New Net MRR}}, se evalúa mediante la siguiente ecuación diferencial de ingresos:
\begin{equation}
    \text{New Net MRR} = \text{New MRR} + \text{Expansion MRR} - \text{Contraction MRR} - \text{Churn MRR}
\end{equation}
cuyos componentes se definen como:
\begin{itemize}
    \item \textbf{\textit{New MRR}:} Ingresos fijos generados por nuevos clubes adquiridos en el mes corriente.
    \item \textbf{\textit{Expansion MRR}:} Crecimiento interno de la cuenta de un club existente. En nuestro modelo, esto ocurre orgánicamente por dos vías: cuando el club crece en su masa societaria o éxito futbolístico (aumentando su abono fijo por escalafón) o, fundamentalmente, por el componente variable: cuando los hinchas incrementan el uso de la plataforma para pagos o reservas, aumentando el volumen del \textit{Fee Transaccional} y del \textit{Marketplace CPA}.
    \item \textbf{\textit{Contraction MRR} (o \textit{Downgrade MRR}):} Pérdida de ingresos por clubes que descienden de categoría o reducen su padrón activo, exigiendo una renegociación a la baja del cargo fijo.
    \item \textbf{\textit{Churn MRR}:} Pérdida total de ingresos recurrentes provocada por clubes que abandonan completamente la plataforma (\textit{Churn}).
\end{itemize}
El objetivo estratégico es garantizar que la suma de ingresos de adquisición y expansión ($\text{New} + \text{Expansion}$) supere holgadamente a la merma por contracción y pérdida de clientes ($\text{Contraction} + \text{Churn}$).

\subsection{Retención institucional (Customer Churn Rate - CCR y Revenue Churn Rate)}

El \textbf{\textit{Customer Churn Rate} (CCR)} mide el porcentaje de clubes que cancelan su suscripción en un período determinado, formalizado mediante la expresión:
\begin{equation}
    \text{CCR} = \frac{\text{Cantidad de clientes perdidos durante el período}}{\text{Cantidad de clientes al inicio del período}}
\end{equation}
A diferencia de las plataformas orientadas al consumidor final (B2C), donde el \textit{churn} suele ser elevado, en el ecosistema B2B institucional de SocioUnido esperamos un CCR extremadamente bajo (inferior al 5\% anual, superando el estándar saludable del mercado SaaS situado entre el 5\% y el 7\% anual).

Esto se debe al concepto de \textbf{``Fricción de Salida'' (\textit{Vendor Lock-in})}: una vez que un club migra su padrón de socios, integra sus finanzas, entrena a su personal de seguridad en el \textit{Smart Access} y educa a sus hinchas a usar el bot de WhatsApp, cambiar a otro proveedor o volver al papel representa un costo operativo e institucional masivo. La plataforma se vuelve el núcleo operativo del club, lo que garantiza una retención natural prolongada.

Adicionalmente, para evitar confusiones entre la pérdida física de cuentas institucionales y la merma económica, incorporamos el seguimiento del \textbf{\textit{Revenue Churn Rate} (o \textit{Gross Dollar Churn})}, que evalúa el impacto financiero directo de la contracción y deserción respecto al MRR del período anterior:
\begin{equation}
    \text{Revenue Churn Rate}_i = \frac{\text{Contraction MRR}_i + \text{Churn MRR}_i}{\text{MRR}_{i-1}}
\end{equation}
donde $\text{MRR}_{i-1}$ representa la facturación recurrente total al cierre del mes inmediatamente anterior. Esta métrica permite auditar si las eventuales pérdidas se concentran en instituciones de gran volumen o en clubes de menor peso contable.

\subsection{Costo de adquisición vs. valor del cliente (CAC y LTV)}

La rentabilidad del modelo de negocio se valida mediante la relación entre cuánto nos cuesta conseguir un club (\textbf{CAC}) y cuánto dinero nos deja ese club durante todo el tiempo que use la plataforma (\textbf{LTV}).

\begin{itemize}
    \item \textbf{\textit{Customer Acquisition Cost} (CAC):} Como se detalla en la estrategia de \textit{Promoción}, la adquisición en SocioUnido no depende de anuncios pagos genéricos (Google Ads), sino de Ventas Consultivas (Outbound), reuniones corporativas, auditorías de recaudación in-situ y alianzas estratégicas. Por lo tanto, el CAC se calcula sumando los costos totales de marketing y ventas del período sobre la cantidad de clientes cerrados:
    \begin{equation}
        \text{CAC} = \frac{\sum \text{Costos de marketing y ventas del período}}{\text{Cantidad de clientes obtenidos en el período}}
    \end{equation}
    En SocioUnido, el numerador incluye los salarios del equipo de ventas y relaciones públicas, viáticos ejecutivos y el costo de oportunidad de bonificar el \textit{Setup Fee} a los \textit{Early Adopters}. Cuanto menor sea el CAC, mayor será el margen bruto del negocio.
    
    \item \textbf{\textit{Lifetime Value} (LTV):} Es el valor monetario total promedio que genera un club a lo largo de toda su vida útil en el ecosistema. Se determina mediante el cociente entre el Ingreso Promedio por Usuario (\textbf{ARPU}) y la tasa de abandono institucional (\textbf{CCR}):
    \begin{equation}
        \text{LTV} = \frac{\text{ARPU}}{\text{CCR}}
    \end{equation}
    donde el \textbf{ARPU} (\textit{Average Monthly Recurring Revenue per User}) es la facturación media por cliente, obtenida del cociente entre el MRR mensual y el total de clubes activos:
    \begin{equation}
        \text{ARPU} = \frac{\text{MRR}}{\text{Cantidad de clientes activos}}
    \end{equation}
    En nuestro modelo B2B, el ARPU de cada club consolida el abono fijo mensual escalonado por padrón, las comisiones del \textit{Fee Transaccional} por cobranzas y los ingresos netos del \textit{Marketplace CPA}.
\end{itemize}

\textbf{Análisis de viabilidad (\textit{LTV/CAC Ratio}):}
Para que SocioUnido sea considerado un negocio altamente rentable y atractivo para inversores, la relación debe cumplir la regla de oro del SaaS:
\begin{equation}
    \text{LTV} > 3 \times \text{CAC}
\end{equation}
asimismo, el período de recupero de la inversión de adquisición (tiempo necesario para compensar el CAC con el MRR neto) debe mantenerse por debajo de los 12 meses ($\text{Meses para recuperar CAC} < 12$), sujeto a la disponibilidad de caja de la empresa. 
Dada la estrategia de comercializar un producto con cobro transaccional recurrente (donde el club aporta dinero continuamente mes a mes) a un sector institucional que rara vez cambia de proveedor (bajando el CCR y extendiendo la vida útil), proyectamos que el índice LTV/CAC de SocioUnido superará ampliamente este umbral, consolidando un modelo de negocio altamente viable y escalable.

\subsection{Velocidad y eficiencia del crecimiento (Quick Ratio)}

A medida que el ecosistema se expanda desde la Fase 1 (Primera B y C) hacia la Fase 2 (Primera Nacional y Liga Profesional), utilizaremos el \textbf{\textit{Quick Ratio}} para medir si estamos creciendo de forma saludable o si simplemente estamos reemplazando clientes perdidos.

La fórmula compara los ingresos que entran versus los que salen:
\begin{equation}
    \text{Quick Ratio} = \frac{\text{New MRR} + \text{Expansion MRR}}{\text{Downgrade MRR} + \text{Churn MRR}}
\end{equation}
donde $\text{Downgrade MRR}$ es equivalente al término $\text{Contraction MRR}$. 

Los rangos de referencia para evaluar la salud de la estructura recurrente son los siguientes:
\begin{itemize}
    \item $\text{Quick Ratio} < 1$: Escenario inviable u obsoleto (\textit{``estamos muertos''}). La merma de ingresos supera a la adquisición y expansión, lo que conduce a una contracción neta que destruye la empresa si perdura en el tiempo.
    \item $1 < \text{Quick Ratio} < 4$: Crecimiento lento o con baja eficiencia. Aunque existe crecimiento neto, la institución o la plataforma requiere mantener altos costos de adquisición para reponer el volumen transaccional perdido.
    \item $\text{Quick Ratio} > 4$: Crecimiento acelerado y de alta eficiencia de retención. Representa el estándar que exige el capital de riesgo (\textit{Venture Capital}) para invertir en empresas SaaS.
\end{itemize}
Nuestra meta corporativa a mediano plazo es sostener un \textbf{Quick Ratio superior a 4}. Lograr este índice (es decir, ingresar \$4 por cada \$1 perdido) validará el éxito de nuestro \textit{Go-To-Market}, demostrando a potenciales grupos inversores que SocioUnido es una empresa que crece a un ritmo acelerado y con alta eficiencia de retención de capital.

\newpage

\subsection{Profit \& Loss (P\&L): Impacto y optimización de costos en las entidades}

El modelo SaaS de SocioUnido no solo se fundamenta en la captación de ingresos recurrentes para la plataforma, sino en una transformación directa sobre el estado de resultados o \textbf{\textit{Profit \& Loss} (P\&L)} de las entidades deportivas clientes. La propuesta de valor actúa simultáneamente sobre la contracción del gasto operativo/fijo y la expansión de la línea de ingresos:

\begin{itemize}
    \item \textbf{Reducción de costos directos, fijos y operativos:}
    \begin{itemize}
        \item \textit{Optimización en sueldos y carga administrativa:} En los clubes del ascenso y ligas metropolitanas, una proporción significativa del presupuesto se destina a horas-hombre ocupadas en tareas repetitivas de cobranza, atención en secretaría y conciliación manual en planillas de cálculo. La implementación del Asistente Conversacional con Procesamiento de Lenguaje Natural (PLN) en WhatsApp para consultas y cobros automáticos, sumado a la automatización contable del Panel E2E, reduce drásticamente la sobrecarga operativa de secretaría y tesorería. Esto permite rediseñar y eficientizar la estructura de gastos en personal, reasignando recursos humanos hacia áreas de mayor valor agregado.
        \item \textit{Ahorro en operativos policiales y de seguridad en accesos:} En los días de partido, la saturación perimetral, las demoras en molinetes y el fraude por entradas o carnets físicos obligan a los clubes a contratar un número elevado de efectivos policiales y personal de seguridad privada para contener cuellos de botella. El módulo de \textbf{Acceso Inteligente (\textit{Smart Access}) con validación TOTP 100\% offline} agiliza el flujo de escaneo a velocidades casi instantáneas sin depender de conectividad ni redes 4G/5G. La eliminación de aglomeraciones y la erradicación del ingreso ilegal por carnet duplicado disminuyen la conflictividad en las puertas, permitiendo a los directivos reducir la dotación requerida y el costo directo del operativo de seguridad por partido.
        \item \textit{Supresión de costos en insumos físicos y servidores locales:} La desmaterialización del carnet de socio en una PWA digital elimina el costo recurrente de impresión de tarjetas PVC, credenciales magnéticas y rollos térmicos. Asimismo, la infraestructura 100\% en la nube ahorra gastos de adquisición, refrigeración y mantenimiento de servidores físicos locales.
    \end{itemize}
    
    \item \textbf{Expansión de ingresos y nuevas vías de monetización:}
    \begin{itemize}
        \item \textit{Recupero de morosidad silenciosa:} El motor predictivo de \textit{Machine Learning} identifica anomalías en la asistencia y el pago antes de que el socio incurra en baja definitiva, activando campañas preventivas de regularización que salvan capital.
        \item \textit{Alquiler digital de espacios institucionales ociosos:} A través de la interfaz del socio, la institución rentabiliza instalaciones compartidas (canchas, salones, quinchos) que permanecían subutilizadas.
        \item \textit{Monetización publicitaria (\textit{Marketplace CPA}):} La integración de un canal comercial y de beneficios segmentado dentro de la aplicación móvil habilita ingresos publicitarios compartidos (10\% para SocioUnido en alianzas directas del club y hasta 75\% para el club en anunciantes provistos por la plataforma), diversificando el ingreso neto corporativo.
    \end{itemize}
\end{itemize}
